% !TEX TS-program = xelatex
% !TEX encoding = UTF-8 Unicode
% !Mode:: "TeX:UTF-8"

\documentclass[11pt]{article}

\setlength{\parindent}{0pt}
\usepackage{xltxtra}
\usepackage{hyperref}
\hypersetup{hidelinks}
\usepackage{url}
\urlstyle{tt}
\usepackage{xcolor}
\definecolor{CVBlue}{RGB}{23,110,191}
\usepackage{calc}
\usepackage{graphicx}
\usepackage{tikz}
\usetikzlibrary{calc}
\usepackage{fontspec}
\usepackage{xeCJK}
\usepackage{enumitem}
\CJKsetecglue{}

%% 主文档字体设置
\setmainfont[
    BoldFont = NotoSerif-Bold.ttf,
]{NotoSerif-Regular.ttf}

% 中文字体设置
\setCJKmainfont[
    BoldFont = NotoSansSC-Bold.ttf,
]{NotoSansSC-Regular.ttf}

\usepackage{relsize}
\usepackage{xspace}

% A4纸，上下左右边距
\usepackage[
    a4paper,
    left=0.95cm,
    right=0.95cm,
    top=1.3cm,
    bottom=0.7cm,
    nohead
]{geometry}

\renewcommand{\baselinestretch}{1.25}

\usepackage{titlesec}
\setlist{noitemsep}
\setlist[itemize]{topsep=0.25em, leftmargin=*}
\setlist[enumerate]{topsep=0.25em, leftmargin=*}

% --- 用于控制【不同项目之间】的垂直距离 ---
\newlength{\interProjectSpacing}
\setlength{\interProjectSpacing}{0.9em}
\newcommand{\projectsep}{\vspace{\interProjectSpacing}}

% --- 用于控制【项目标题】与下方【项目描述】的距离 ---
\newlength{\intraProjectTitleSep}
\setlength{\intraProjectTitleSep}{0.4em}
\newcommand{\titlebreak}{\\[\intraProjectTitleSep]}

% --- 用于控制【项目描述】与下方【要点列表】的距离 ---
\newlength{\intraProjectListTopSep}
\setlength{\intraProjectListTopSep}{0.2em}

% =======================================================================

\titleformat{\section}
{\large\bfseries\raggedright}
{}{0em}
{}
[{\color{CVBlue}\titlerule}]
\titlespacing*{\section}{0cm}{*1.6}{*1.2}

\begin{document}
\pagenumbering{gobble}

%%%% 利用tikz来定位照片（右上角）
\begin{tikzpicture}[remember picture, overlay]
    \node[anchor = north east] at ($(current page.north east)+(-1.2cm,-1.5cm)$) {\includegraphics[height=3cm]{avatar.jpg}};
  \end{tikzpicture}%
  %%%% 利用tikz来定位学校Logo，这里只在第一页显示（左上角）
  \begin{tikzpicture}[remember picture, overlay]
    \node[anchor = north west] at ($(current page.north west)+(0.3cm,+0.9cm)$) {\includegraphics[height=6cm]{zju.png}};
  \end{tikzpicture}%

\centerline{\LARGE\bfseries{张平安}}
\centerline{\normalsize{\textsection\,Phone~177-8673-5893 \quad \textbullet\,Email~\href{mailto:2293450367@qq.com}{2293450367@qq.com}}}
\centerline{\normalsize{\textbullet\, 浙江大学电子信息工程硕士 \quad \textbullet\, 浙江·杭州}}

\section{\color{CVBlue}\ensuremath{\blacksquare}\ 教育背景}

\textbf{浙江大学} \hfill 2024.09 -- 2027.06 \titlebreak
电子信息工程 \quad \textit{工学硕士}
\begin{itemize}[nosep, topsep=\intraProjectListTopSep]
    \item \textbf{荣誉奖励}：浙江大学硕士一等学业奖学金（专业前 40\%）
\end{itemize}

\projectsep

\textbf{武汉科技大学} \hfill 2020.09 -- 2024.06 \titlebreak
电子信息工程 \quad \textit{工学学士}
\begin{itemize}[nosep, topsep=\intraProjectListTopSep]
    \item \textbf{学业成绩}：专业综合排名前 10\%，连续获得国家励志奖学金、校一等奖学金
\end{itemize}

\section{\color{CVBlue}\ensuremath{\blacksquare}\ 项目与实习经历}
\textbf{循澜科技有限公司 | 基于ROS的AGV自主巡检与自动回充系统} \hfill 2026.03 -- 2026.06 \titlebreak
\textit{算法开发}\titlebreak
\begin{itemize}[nosep, topsep=\intraProjectListTopSep]
    \item 基于激光雷达与 IMU 部署 Omni-LIVO 建图流程，经 ROS bag 录制回放生成点云地图并转换为二维栅格地图，接入 AMCL 定位链路。
    \item 针对转弯后定位漂移与位姿跳变，调节 AMCL 粒子数与里程计噪声；调优 TEB 速度与障碍物代价，消除轨迹贴墙与频繁急停。
    \item 经多轮实车测试成功对准充电桩并触发充电，偏航角误差控制在 2° 以内、横向误差控制在 2 cm 以内。
    \item 接入红外热成像相机并配置静态 TF，以 80℃ 阈值加 3 帧确认实现 15 fps 火情监测；打通导航、底盘与 MQTT 链路。
\end{itemize}

\projectsep

\textbf{浙江大学金华研究院 | 穿戴式生理声学时序分析与感知平台} \hfill 2024.11 -- 2026.01 \titlebreak
\textit{算法与系统开发}\titlebreak
\begin{itemize}[nosep, topsep=\intraProjectListTopSep]
    \item 将 8 kHz/16 bit 生理声学信号切分为 1 s 样本，完成滤波与归一化，提取 24 维 MFCC 及差分，构建 3×24×512 特征图。
    \item 按受试者维度以 6:2:2 划分训练/验证/测试集，以 PyTorch 构建 ResNet-LSTM 二分类模型，引入 Coordinate Attention 与 Focal Loss。
    \item 采用加噪、时移及频率遮挡增强，测试集 F1 与 AUC 达 0.93--0.95。
    \item 基于 nRF52840 与 Zephyr RTOS 开发 PDM 采集模块，以 DMA 双缓冲降低 CPU 开销。
    \item 静态待机电流降至 15 μA、通信功耗降低 60%；完成 INT8 量化及 FPGA 推理部署验证。
\end{itemize}

\section{\color{CVBlue}\ensuremath{\blacksquare}\ 技术能力}
编程语言：Python、PyTorch、C（熟练），C++（基础语法）；机器人导航：ROS Topic / 参数服务器 / TF / ROS bag，Omni-LIVO、AMCL、TEB、AprilTag、阿克曼底盘；信号处理：音频切分、滤波去噪、MFCC 与差分、小波维纳滤波、Focal Loss；嵌入式与部署：nRF52840、Zephyr RTOS、PDM、DMA、BLE、INT8 量化与 FPGA 部署；工具链：Git、Docker、Gazebo、MQTT。

\section{\color{CVBlue}\ensuremath{\blacksquare}\ 荣誉奖励与语言能力}
\begin{itemize}[nosep]
    \item “大唐杯”5G 大赛省二等奖、蓝桥杯省三等奖、全国大学生英语竞赛（NECCS）全国三等奖；CET-6 580 分。
\end{itemize}

\end{document}
